%%%%%%%%%%%%%%%%%%%%%%%%%%%%%%%%%%%%%%%%%%%%%%%%%%%%%%%%%%%%%%%%%%%%%%%%%%%%%%%
% CHAPTER 6: THE REFERENCE STRUCTURE C*
%%%%%%%%%%%%%%%%%%%%%%%%%%%%%%%%%%%%%%%%%%%%%%%%%%%%%%%%%%%%%%%%%%%%%%%%%%%%%%%
% Version: 2.0 (January 2026)
%
% Purpose: Conceptual introduction to C* as the self-consistent reference
%          structure. All formal content (axioms, theorems, proofs) is in
%          Appendix D (FORMAL-CSTAR).
%
% Primary Appendix: D - FORMAL-CSTAR (Complete C* Theory)
% Secondary Appendices: B (Complementarity), BBB (Calibration)
%
% Prerequisites: Ch. 5 (Complementarity)
% Leads to: Ch. 7 (Coherence Index), Ch. 9 (Context)
%
% Chapter Type: B (Foundation Chapter - conceptual overview)
%%%%%%%%%%%%%%%%%%%%%%%%%%%%%%%%%%%%%%%%%%%%%%%%%%%%%%%%%%%%%%%%%%%%%%%%%%%%%%%

\chapter{The Reference Structure $C^*$}
\label{ch:reference-structure}

%==============================================================================
% DIE 9 FRAGEN NAVIGATION BOX
%==============================================================================
\BCMQuestionNav{6}

%==============================================================================
% 9C CORE CONNECTION BOX
%==============================================================================
\begin{tcolorbox}[
    colback=blue!5!white,
    colframe=blue!75!black,
    title={\textbf{9C CORE Connection: Foundation for HOW}},
    fonttitle=\bfseries
]
\textbf{Foundation Question:} What is the reference structure $C^*$?

\medskip
\begin{tabular}{ll}
\textbf{CORE Appendix:} & D (FORMAL-CSTAR) \\
\textbf{Output:} & Self-consistent structure $C^* = \Phi(C^*, \Psi)$ \\
\textbf{Key Insight:} & $C^*$ is derived, not assumed \\
\textbf{Novel Contribution:} & Three convergent derivations
\end{tabular}

\medskip
\textbf{The Core Equation:}
\begin{center}
$C^* = \Phi(C^*, \Psi)$ \quad (Fixed-point: beliefs reproduce themselves)
\end{center}
\end{tcolorbox}

%==============================================================================
% APPENDIX REFERENCES BOX
%==============================================================================
\begin{tcolorbox}[
    colback=yellow!10!white,
    colframe=orange!75!black,
    title={\textbf{Appendix References for This Chapter}},
    fonttitle=\bfseries
]
\begin{tabular}{lll}
\textbf{Appendix} & \textbf{Content} & \textbf{Section} \\
\midrule
D (FORMAL-CSTAR) & Complete C* theory: axioms, theorems, proofs & All \\
B (CORE-HOW) & Complementarity structure $\Gamma$ & \S\ref{sec:cstar-derivations} \\
BBB (CORE-WHERE) & Calibration of $C^*(\Psi)$ & \S\ref{sec:cstar-calibration}
\end{tabular}

\medskip
\textbf{Reading Guide:} This chapter provides intuition; Appendix D contains
the complete mathematical theory with all proofs.
\end{tcolorbox}


%%%%%%%%%%%%%%%%%%%%%%%%%%%%%%%%%%%%%%%%%%%%%%%%%%%%%%%%%%%%%%%%%%%%%%%%%%%%%%%
%                    PART I: INTRODUCTION
%%%%%%%%%%%%%%%%%%%%%%%%%%%%%%%%%%%%%%%%%%%%%%%%%%%%%%%%%%%%%%%%%%%%%%%%%%%%%%%

\section{Why $C^*$ Matters}
\label{sec:cstar-intro}

The coherence index $K = 1 - \|C - C^*\|/\|C^*\|$ measures deviation from a
reference structure $C^*$. For this to be scientifically meaningful, $C^*$
must be \emph{derived}, not assumed.

\begin{intuition}[The Trust Example]
Merchants deciding whether to extend credit:
\begin{itemize}[nosep]
    \item \textbf{Beliefs match reality:} Trust is rewarded---beliefs reproduce
    \item \textbf{Too optimistic:} Get cheated, revise downward
    \item \textbf{Too pessimistic:} Miss opportunities, discover cooperation works
\end{itemize}
In all cases, beliefs gravitate toward $C^*$---the structure that reproduces itself.
\end{intuition}

\paragraph{The Stigler-Becker Parallel.}
\citet{stiglerbecker1977} argued that economists should assume stable preferences.
We preserve this insight at a higher level: not stable preferences $U$, but stable
reference structure $C^*(\Psi)$. The key advance: $C^*$ is not exogenously assumed
but derived as the unique self-consistent configuration.


%%%%%%%%%%%%%%%%%%%%%%%%%%%%%%%%%%%%%%%%%%%%%%%%%%%%%%%%%%%%%%%%%%%%%%%%%%%%%%%
%                    PART II: THREE DERIVATIONS
%%%%%%%%%%%%%%%%%%%%%%%%%%%%%%%%%%%%%%%%%%%%%%%%%%%%%%%%%%%%%%%%%%%%%%%%%%%%%%%

\section{Three Derivations of $C^*$}
\label{sec:cstar-derivations}

Appendix D provides three independent derivations that all yield the same $C^*$:

\paragraph{1. Axiomatic Derivation (Self-Consistency).}
$C^*$ is the unique fixed point: $C^* = \Phi(C^*, \Psi)$. Under regularity,
reciprocity, and contraction axioms (D-REG, D-REC, D-CON), the Banach
fixed-point theorem guarantees existence and uniqueness.
\emph{See Appendix D, Part II.}

\paragraph{2. Evolutionary Derivation (Stability).}
$C^*$ is the evolutionarily stable structure (ESS). No mutant belief $C' \neq C^*$
can invade a population at $C^*$. Naive cooperators get exploited; paranoid
defectors miss cooperative gains. Both are outcompeted by accurate beliefs.
\emph{See Appendix D, Part III.}

\paragraph{3. Empirical Derivation (Long-Run Average).}
$C^*$ is the ergodic limit: $C^*_{emp} = \lim_{T\to\infty} \frac{1}{T}\sum_t C_t$.
Under adaptive dynamics with unique stationary distribution, the long-run
average converges to the theoretical $C^*$.
\emph{See Appendix D, Part IV.}

\begin{theorem}[Convergence---Appendix D, Theorem D.5]
Under the foundational axioms and ergodicity:
$C^*_{axiomatic} = C^*_{evolutionary} = C^*_{empirical}$.
\end{theorem}


%%%%%%%%%%%%%%%%%%%%%%%%%%%%%%%%%%%%%%%%%%%%%%%%%%%%%%%%%%%%%%%%%%%%%%%%%%%%%%%
%                    PART III: CALIBRATION LIMITS
%%%%%%%%%%%%%%%%%%%%%%%%%%%%%%%%%%%%%%%%%%%%%%%%%%%%%%%%%%%%%%%%%%%%%%%%%%%%%%%

\section{Calibration Requirements}
\label{sec:cstar-calibration}

$C^*$ cannot be learned from text alone. Each derivation requires empirical input:
\begin{itemize}[nosep]
    \item Axiomatic: Parameters $\alpha, \beta, \gamma$ (magnitudes not in axioms)
    \item Evolutionary: Fitness landscape (selection pressures not documented)
    \item Empirical: Behavioral observations (effect sizes scattered)
\end{itemize}

\emph{See Appendix BBB (CORE-WHERE) for calibration; Appendix AN for LLM-MC estimation.}


%%%%%%%%%%%%%%%%%%%%%%%%%%%%%%%%%%%%%%%%%%%%%%%%%%%%%%%%%%%%%%%%%%%%%%%%%%%%%%%
%                    PART IV: 9C INTEGRATION
%%%%%%%%%%%%%%%%%%%%%%%%%%%%%%%%%%%%%%%%%%%%%%%%%%%%%%%%%%%%%%%%%%%%%%%%%%%%%%%

\section{9C Integration}
\label{sec:cstar-integration}

\begin{table}[htbp]
\centering
\caption{Chapter 6 Integration with 9C CORE Structure}
\label{tab:ch6-8c-integration}
\small
\begin{tabular}{@{}llp{5cm}l@{}}
\toprule
\textbf{CORE} & \textbf{Question} & \textbf{Connection} & \textbf{Ref} \\
\midrule
\cellcolor{purple!20}WHO & Who has utility? & Agents in $C^*$ computation & AAA \\
\cellcolor{red!20}WHAT & What is utility? & Dimensions in $C^*$ matrix & C \\
\cellcolor{blue!20}HOW & How interact? & $C^*$ is the stable $\Gamma$-matrix & B \\
\cellcolor{green!20}WHEN & Context? & $C^* = C^*(\Psi)$ depends on context & V \\
\cellcolor{gray!20}WHERE & Parameters? & Calibration of $C^*$ & BBB \\
\cellcolor{orange!20}AWARE & Awareness? & Beliefs about $C^*$ matter & AU \\
\cellcolor{violet!20}READY & Action threshold? & $C^*$ determines equilibrium & AV \\
\cellcolor{teal!20}STAGE & Journey phase? & $C^*$ defines stable configs & AW \\
\bottomrule
\end{tabular}
\end{table}


%%%%%%%%%%%%%%%%%%%%%%%%%%%%%%%%%%%%%%%%%%%%%%%%%%%%%%%%%%%%%%%%%%%%%%%%%%%%%%%
%                    PART V: SUMMARY
%%%%%%%%%%%%%%%%%%%%%%%%%%%%%%%%%%%%%%%%%%%%%%%%%%%%%%%%%%%%%%%%%%%%%%%%%%%%%%%

\section{Summary}
\label{sec:cstar-summary}

\begin{tcolorbox}[
    colback=blue!5!white,
    colframe=blue!75!black,
    title={\textbf{Chapter 6 Summary}}
]
\textbf{Core Insight:} $C^*$ is the unique self-consistent reference structure.

\textbf{Three Derivations} (all in Appendix D):
\begin{enumerate}[nosep]
    \item Axiomatic: Fixed-point $C^* = \Phi(C^*, \Psi)$
    \item Evolutionary: ESS condition
    \item Empirical: Ergodic long-run average
\end{enumerate}

\textbf{Key Results:}
\begin{itemize}[nosep]
    \item $C^*$ exists, is unique, and is context-dependent
    \item Three derivations converge on the same $C^*$
    \item $C^*$ requires calibration---cannot be learned from text
\end{itemize}

\textbf{Transition:} Chapter 7 uses $C^*$ to define coherence dynamics $K(t)$.
\end{tcolorbox}

%------------------------------------------------------------------------------
% READING PATH BOX
%------------------------------------------------------------------------------
\begin{tcolorbox}[
    colback=gray!10!white,
    colframe=gray!75!black,
    title={\textbf{Reading Paths}}
]
\textbf{For Complete Theory:} $\rightarrow$ Appendix D (FORMAL-CSTAR)

\textbf{For Practitioners:} $\rightarrow$ Chapter 7 (coherence) $\rightarrow$ Chapter 9 (context)

\textbf{For Calibration:} $\rightarrow$ Appendix BBB $\rightarrow$ Appendix AN
\end{tcolorbox}

%%%%%%%%%%%%%%%%%%%%%%%%%%%%%%%%%%%%%%%%%%%%%%%%%%%%%%%%%%%%%%%%%%%%%%%%%%%%%%%
